% Delta-Neutral Trading Strategy Using Binomial Option Pricing
% Author: Roan McReynolds, Jayden Thinh-To, Jordan Odorico, Charles Seitz
% Date: March 2026
%
% FIGURE DIRECTORIES (Overleaf):
%   figures/          <- copy all files from outputs/reports/figures/
%   convergence/      <- copy all files from outputs/convergence/
%   param_sweep/      <- copy heatmap_sharpe.png, heatmap_maxdd.png from outputs/param_sweep/

\documentclass[10pt,twocolumn]{article}
\usepackage[margin=1in,top=1in,columnsep=0.25in]{geometry}
\usepackage{amsmath}
\usepackage{amssymb}
\usepackage{graphicx}
\usepackage{float}
\usepackage{booktabs}
\usepackage{hyperref}
\usepackage{enumitem}
\usepackage{listings}
\usepackage{xcolor}
\usepackage{caption}
\usepackage{subcaption}
\usepackage{titlesec}
\usepackage{abstract}
\usepackage{fancyhdr}
\usepackage{microtype}
\usepackage{authblk}
\usepackage{mathtools}
\usepackage{cuted}
\usepackage{multirow}

\hypersetup{
    colorlinks=true,
    linkcolor=blue,
    citecolor=blue,
    urlcolor=blue
}

% Make \begin{thebibliography} produce a numbered section heading
\makeatletter
\renewenvironment{thebibliography}[1]{%
  \section{\refname}%
  \@mkboth{\MakeUppercase\refname}{\MakeUppercase\refname}%
  \list{\@biblabel{\@arabic\c@enumiv}}{%
    \settowidth\labelwidth{\@biblabel{#1}}%
    \leftmargin\labelwidth
    \advance\leftmargin\labelsep
    \@openbib@code
    \usecounter{enumiv}%
    \let\p@enumiv\@empty
    \renewcommand\theenumiv{\@arabic\c@enumiv}}%
  \sloppy
  \clubpenalty4000
  \@clubpenalty\clubpenalty
  \widowpenalty4000%
  \sfcode`\.\@m}{%
  \def\@noitemerr{\@latex@warning{Empty `thebibliography' environment}}%
  \endlist}
\makeatother

\titleformat{\section}{\normalfont\bfseries\centering}{}{0em}{\thesection.~}
\titleformat{\subsection}{\normalfont\bfseries}{\thesubsection}{1em}{}
\titleformat{\subsubsection}{\normalfont\itshape}{\thesubsubsection}{1em}{}

\titlespacing*{\section}{0pt}{8pt}{4pt}
\titlespacing*{\subsection}{0pt}{6pt}{2pt}
\titlespacing*{\subsubsection}{0pt}{4pt}{2pt}

\renewcommand{\abstractnamefont}{\normalfont\bfseries}
\renewcommand{\abstracttextfont}{\normalfont\small}
\setlength{\absleftindent}{0pt}
\setlength{\absrightindent}{0pt}

\pagestyle{fancy}
\fancyhf{}
\fancyhead[L]{\small\textit{Delta-Neutral Trading Strategy}}
\fancyhead[R]{\small QUANTT Research}
\fancyfoot[C]{\small\thepage}
\renewcommand{\headrulewidth}{0.4pt}

\lstset{
    language=Python,
    basicstyle=\ttfamily\footnotesize,
    keywordstyle=\color{blue},
    commentstyle=\color{gray},
    stringstyle=\color{red},
    numbers=left,
    numberstyle=\tiny\color{gray},
    stepnumber=1,
    frame=single,
    breaklines=true,
    captionpos=b
}

\setlength{\parindent}{0pt}
\setlength{\parskip}{3pt}

\setlength{\floatsep}{6pt plus 2pt minus 2pt}
\setlength{\textfloatsep}{8pt plus 2pt minus 4pt}
\setlength{\intextsep}{6pt plus 2pt minus 2pt}
\setlength{\dblfloatsep}{6pt plus 2pt minus 2pt}
\setlength{\dbltextfloatsep}{8pt plus 2pt minus 4pt}
\captionsetup{font=small, skip=3pt}

% Two figure search paths: main report figures and convergence figures
\graphicspath{{figures/}{convergence/}{param_sweep/}}


\title{
    \large\bfseries
    From Price Mispricing to Volatility Surface Arbitrage: A CRR Binomial
    Framework for Delta-Neutral Options Research
}

\author[1]{Roan McReynolds}
\author[1]{Jayden Thinh-To}
\author[1]{Jordan Odorico}
\author[1]{Charles Seitz}
\affil[1]{QUANTT --- Queen's University Algorithmic \& Network Trading Team}

\date{March 2026}

% ══════════════════════════════════════════════════════════════════════════════
\begin{document}
% ══════════════════════════════════════════════════════════════════════════════

\twocolumn[
  \maketitle
  \thispagestyle{fancy}

  \begin{onecolabstract}
We develop and evaluate a progression of delta-neutral options strategies
based on the Cox-Ross-Rubinstein (CRR) binomial model.  A baseline Z-score
mispricing strategy, entering when market price deviates from a
rolling historical-volatility CRR fair value, produces no persistent edge
on SPY options during 2023--2024, with mispricings structurally attributable
to the volatility risk premium rather than model arbitrage.  We redesign
around a same-day cross-sectional implied volatility surface: implied
volatilities are solved via CRR bisection across the full option chain, a
polynomial surface is fitted in log-moneyness and time-to-expiry space, and
contracts deviating materially from surface-fitted IV are traded against the
surface-fair CRR price with a CRR-delta hedge.  This eliminates the
volatility risk premium bias inherent in the historical-volatility baseline.

Strategy hyperparameters are selected via a grid search on a 2022 training
window (a bear-market year, SPY $-18\%$) subject to a maximum-drawdown
constraint, producing robust parameters that are subsequently validated
out-of-sample on 2023--2024.  Over the full three-year period 2022--2024
(753 trading days encompassing three qualitatively distinct market regimes),
the primary PUT\_30\_60\_OTM bucket accumulates \$126,523 (annualized Sharpe
3.41, Sortino 3.96, Calmar 8.16), with out-of-sample Sharpe of 2.03 on the
unseen 2023--2024 period.  The CALL\_120\_150\_ATM bucket accumulates
\$47,488 (Sharpe 1.49).  Both buckets substantially outperform a passive SPY
buy-and-hold on a risk-adjusted basis.  All results assume mid-price
execution with a 25\,bps synthetic half-spread and no brokerage commissions;
a transaction-cost sensitivity analysis in Section~6.7 quantifies the impact
of realistic frictions.

  \vspace{6pt}
  \noindent\rule{\linewidth}{0.4pt}
  \vspace{4pt}
  \end{onecolabstract}
]


% ══════════════════════════════════════════════════════════════════════════════
\section{Introduction to the Binomial Option Pricing Model}
% ══════════════════════════════════════════════════════════════════════════════

The binomial option pricing model provides a discrete-time framework for
valuing derivatives and constructing dynamic hedging strategies. This section
presents the mathematical foundations of the CRR parameterization, its
implementation structure, and the calculation of the Greeks used throughout
both strategy generations described in this paper.

\subsection{Model Overview and Motivation}

The binomial model, introduced by Cox, Ross, and Rubinstein \cite{crr1979},
discretizes asset price evolution into a recombining lattice of up/down moves,
recovering option values by risk-neutral backward induction. Its tree structure natively accommodates American early-exercise features that Black-Scholes cannot
price directly \cite{hull2018}, and as $N \to \infty$ it converges in
distribution to geometric Brownian motion \cite{black1973}. The CRR pricer
also serves as the numerical engine for inverting market prices to implied
volatilities via bisection, a computation that underpins both strategy
generations discussed below.

\subsection{Mathematical Framework}

\subsubsection{Underlying Asset Price Tree}

The underlying asset price evolves according to a multiplicative binomial
process. Given an initial price $S_0$, the price at each subsequent node is
determined by the application of up and down factors under the CRR
parameterization.

\textbf{CRR Parameterization.} The up and down factors are calibrated to
match the annualized volatility $\sigma$ of the underlying:
%
\begin{equation}
u = e^{\sigma \sqrt{\Delta t}}, \qquad d = e^{-\sigma\sqrt{\Delta t}} =
\frac{1}{u},
\end{equation}
%
where $\Delta t = T/N$, $T$ is time to expiration, and $N$ is the number of
steps. At time step $i$ and state $j$ (the number of down moves), the
underlying price is:
%
\begin{equation}
S_{i,j} = S_0 \cdot u^{i-j} \cdot d^{j}.
\end{equation}

\textbf{Recombining Tree Property.} Because $u \cdot d = 1$, the tree
recombines: an up move followed by a down move yields the same node as a
down move followed by an up move. This reduces computational complexity
from $O(2^N)$ to $O(N^2)$ nodes, making large step counts tractable.

\subsubsection{Risk-Neutral Probability}

Under the risk-neutral measure, the expected return of the underlying asset
equals the cost of carry $r - q$, where $r$ is the risk-free rate and $q$ is
the continuous dividend yield. The no-arbitrage condition
%
\begin{equation}
e^{(r-q)\Delta t} S_{i,j} = p\,S_{i,j}\cdot u + (1-p)\,S_{i,j}\cdot d
\end{equation}
%
yields the risk-neutral up-probability:
%
\begin{equation}
p = \frac{e^{(r-q)\Delta t} - d}{u - d}.
\end{equation}
%
This probability is a pricing device constructed to be consistent with the
absence of arbitrage; it does not represent the physical probability of an
upward move \cite{hull2018}.

\subsubsection{Option Payoff at Expiration}

At the terminal nodes $i = N$, the option value is its intrinsic payoff.
For a call:
%
\begin{equation}
V_{N,j} = \max(S_{N,j} - K,\;0),
\end{equation}
%
and for a put:
%
\begin{equation}
V_{N,j} = \max(K - S_{N,j},\;0),
\end{equation}
%
where $K$ is the strike price.

\subsection{Backward Induction for Option Valuation}

Option values at earlier nodes are recovered by working backwards from
expiration using discounted risk-neutral expectations.

\subsubsection{European Options}

For European-style options, exercisable only at expiration, the value at
node $(i,j)$ is:
%
\begin{equation}
V_{i,j} = e^{-r\Delta t}\bigl[p\cdot V_{i+1,j}
          + (1-p)\cdot V_{i+1,j+1}\bigr].
\end{equation}

\subsubsection{American Options}

For American-style options the holder maximizes value by comparing immediate
exercise against continuation:
%
\begin{equation}
V_{i,j} = \max\bigl(\text{Exercise Value},\;\text{Hold Value}\bigr),
\end{equation}
%
where the hold value is the same discounted risk-neutral expectation as in
the European case, and the exercise value is $\max(S_{i,j}-K,0)$ for a call
or $\max(K-S_{i,j},0)$ for a put. The option to exercise early is always
weakly valuable, so American options price weakly above their European
counterparts. All SPY options analyzed in this paper are American-style.

\subsection{Discrete Delta Calculation}

The delta $\Delta$ measures the sensitivity of the option price to the
underlying, defined continuously as $\partial V / \partial S$. In the
binomial lattice it is computed as a finite-difference approximation at the
root node across the two first-step successor states:
%
\begin{equation}
\Delta = \frac{V_{1,0} - V_{1,1}}{S_{1,0} - S_{1,1}}.
\end{equation}
%
Delta ranges in $[0,1]$ for calls and $[-1,0]$ for puts. Deep in-the-money
options have $|\Delta| \approx 1$; far out-of-the-money options have
$|\Delta| \approx 0$; at-the-money options cluster near $|\Delta| = 0.5$.
The hedge share count per contract is:
%
\begin{equation}
\text{Shares} = -\Delta \times 100,
\end{equation}
%
where the negative sign reflects the requirement to take the opposite
directional position in the underlying in order to offset the option's
exposure.

\subsection{Implied Volatility via CRR Bisection}

A critical operation used in both strategy generations is inverting the CRR
pricing function to recover implied volatility (IV) from an observed market
price. Because the CRR pricer has no closed-form inverse in $\sigma$, we
employ numerical bisection. For a given market price $P^{\text{mkt}}$ and all
other parameters fixed, we seek $\hat{\sigma}$ satisfying:
%
\begin{equation}
\text{CRR}(S_0, K, T, r, q, \hat{\sigma}) = P^{\text{mkt}}.
\end{equation}
%
The bisection exploits the fact that the CRR price is strictly monotone
increasing in $\sigma$: for any two candidates $\sigma_a < \sigma_b$,
$\text{CRR}(\sigma_a) < \text{CRR}(\sigma_b)$. An initial bracket
$[\sigma_{\text{lo}}, \sigma_{\text{hi}}] = [0.01, 2.50]$ is established,
and at each iteration the midpoint $\sigma_m = (\sigma_{\text{lo}} +
\sigma_{\text{hi}})/2$ is evaluated. If $\text{CRR}(\sigma_m) >
P^{\text{mkt}}$, the upper bracket is tightened to $\sigma_m$; otherwise
the lower bracket is raised to $\sigma_m$. After at most 100 iterations,
the bracket width is reduced below a tolerance of $10^{-5}$, yielding
$\hat{\sigma}$ to the precision required for surface fitting and signal
generation. This procedure is the backbone of both the time-series IV
z-score strategy and the cross-sectional surface relative-value strategy.

\subsection{Implementation Structure and Performance}

The Python implementation uses Numba \cite{lam2015} just-in-time compilation
on the core tree construction and backward induction routines. Stock and
option trees are allocated as $(N+1) \times (N+1)$ arrays and filled in
vectorized loops that compile to native machine code on first call. On
subsequent calls the compiled kernel is cached, making the per-contract
pricing cost negligible relative to the cost of data retrieval. For the IV
bisection, each call to \texttt{implied\_vol\_crr} invokes the compiled pricer
up to 100 times; in practice 30--50 iterations suffice for the $10^{-5}$
tolerance.

\subsection{Model Assumptions and Limitations}

The CRR model assumes constant volatility, a fixed risk-free rate and dividend
yield, no transaction costs, and continuous asset divisibility. Realized
volatility is stochastic and mean-reverting \cite{andersen2003,bollerslev1986},
and trading frictions are non-trivial in practice \cite{whalley1997}. Most
importantly, a single scalar volatility input cannot reproduce the cross-sectional
structure of market option prices, the empirical volatility smile and term
structure, motivating the surface-based strategy in
Section~\ref{sec:surface_strategy}.


% ══════════════════════════════════════════════════════════════════════════════
\section{Delta-Neutral Hedging}
% ══════════════════════════════════════════════════════════════════════════════

\subsection{Concept of Market Neutrality}

A delta-neutral portfolio is constructed to be instantaneously insensitive
to small movements in the underlying asset. By holding a position in an
option alongside an offsetting position in the underlying sized by the
option's delta, a portfolio manager can isolate exposure to other
risk factors, volatility, time decay, or cross-sectional pricing
anomalies, while eliminating first-order directional risk
\cite{hull2018}. The strategy returns are thus driven by the accuracy
of the non-directional signal rather than by the level or direction of
the underlying, a property essential for rigorous evaluation of any
mispricing-based approach.

\subsection{Constructing the Delta-Neutral Portfolio}

Upon entering a position in $C$ contracts on one side of an option, an
offsetting stock position is established. For a long call position with
delta $\Delta > 0$, the delta-neutral hedge requires shorting the
underlying:
%
\begin{equation}
\text{HedgeShares} = -\Delta \times 100 \times C.
\end{equation}
%
For a short option position of side $s \in \{-1, +1\}$ where $s = -1$
denotes a short and $s = +1$ denotes a long:
%
\begin{equation}
\text{HedgeShares} = -s \times \Delta \times 100 \times C.
\end{equation}
%
The combined portfolio value at initiation is approximately invariant to
infinitesimal moves in $S$, satisfying the neutrality condition
$\partial \Pi / \partial S \approx 0$.

\subsection{Dynamic Rebalancing}

Delta is not constant. As $S$ moves, $\Delta$ changes at a rate governed by
the option's gamma $\Gamma = \partial^2 V / \partial S^2$. Continuous
rebalancing to maintain exact delta-neutrality is theoretically optimal but
practically infeasible due to transaction costs \cite{whalley1997}. In our
implementation, the hedge is rebalanced discretely whenever delta drifts
beyond a threshold from its last recorded value:
%
\begin{equation}
|\Delta_t - \Delta_{\text{last}}| > \delta_{\text{threshold}}.
\end{equation}
%
At each rebalance event, the stock position is adjusted by the difference in
required shares:
%
\begin{equation}
\Delta\text{Shares} = \bigl(-s \times C \times 100 \times \Delta_t\bigr)
                    - \text{HedgeShares}_{\text{last}},
\end{equation}
%
and $\Delta_{\text{last}}$ is updated to $\Delta_t$. The rebalancing
threshold trades off hedging precision against transaction cost accumulation;
its calibration is discussed separately in each strategy section below.
The cumulative impact of discrete rebalancing on net P\&L across a range of
assumed spread levels is quantified in Section~6.7.


% ══════════════════════════════════════════════════════════════════════════════
\section{Strategy I: Z-Score Price Mispricing}
\label{sec:zscore_strategy}
% ══════════════════════════════════════════════════════════════════════════════

\subsection{Motivation}

The first strategy attempts to exploit mean-reversion in the spread between
an option's market price and the fair value produced by the CRR model when
supplied with rolling historical volatility. The intuition is that if the
market price deviates unusually far from the model's fair value in a
statistically normalized sense, it is likely to revert, and a delta-hedged
position can capture that convergence \cite{stein1989}.

\subsection{Signal Construction}

\subsubsection{Model-Implied Fair Value}

On each trading day $t$, for each contract under evaluation, the CRR model
is priced using the 30-day rolling historical volatility estimate
$\hat{\sigma}_t^{\text{HV}}$ computed from underlying log-returns:
%
\begin{equation}
\hat{\sigma}_t^{\text{HV}} = \sqrt{252} \cdot
\text{StdDev}\!\left(r_{t-29}, \ldots, r_t\right), \quad
r_t = \ln\!\left(\frac{S_t}{S_{t-1}}\right).
\end{equation}
%
The model-implied fair value is then:
%
\begin{equation}
P_t^{\text{model}} = \text{CRR}\!\left(S_t, K, T_t, r,
                     q, \hat{\sigma}_t^{\text{HV}}\right),
\end{equation}
%
where $T_t = (\text{expiry} - t) / 252$ is updated daily.

\subsubsection{Raw Mispricing and Z-Score Normalization}

The raw mispricing is defined as:
%
\begin{equation}
\epsilon_t = P_t^{\text{market}} - P_t^{\text{model}}.
\end{equation}
%
Over a rolling lookback window of $L$ observations, the mean and standard
deviation of past mispricings are computed:
%
\begin{equation}
\mu_\epsilon = \frac{1}{L}\sum_{i=1}^{L}\epsilon_{t-i}, \qquad
\sigma_\epsilon = \sqrt{\frac{1}{L}\sum_{i=1}^{L}
                 (\epsilon_{t-i}-\mu_\epsilon)^2}.
\end{equation}
%
The z-score of the current mispricing is then:
%
\begin{equation}
z_t = \frac{\epsilon_t - \mu_\epsilon}{\sigma_\epsilon}.
\end{equation}

\subsection{Entry and Exit Rules}

A position is entered when $|z_t|$ crosses an entry threshold
$z_{\text{entry}}$:
\begin{itemize}[noitemsep]
  \item $z_t > z_{\text{entry}}$: option is statistically overpriced
        $\Rightarrow$ sell the option, buy $|\text{HedgeShares}|$
        of the underlying.
  \item $z_t < -z_{\text{entry}}$: option is statistically underpriced
        $\Rightarrow$ buy the option, short $|\text{HedgeShares}|$
        of the underlying.
\end{itemize}
The position is closed when $|z_t| < z_{\text{exit}}$, indicating the
mispricing has reverted sufficiently. Default thresholds are
$z_{\text{entry}} = 2.0$ and $z_{\text{exit}} = 0.5$.

\subsection{Portfolio Accounting}

The mark-to-market equity evolution for an open position is:
%
\begin{equation}
\Pi_t =
\begin{cases}
\text{Proceeds} - 100C\cdot P_t^{\text{mkt}} + H \cdot S_t,
  & \text{short option},\\[4pt]
100C\cdot P_t^{\text{mkt}} - \text{Cost} + H \cdot S_t,
  & \text{long option},
\end{cases}
\end{equation}
%
where $C$ is contracts, $H$ is hedge shares, and Proceeds (Cost) is the
initial cash flow from selling (purchasing) the option. Daily P\&L is
the first difference of $\Pi_t$, and the cumulative equity curve is its
running sum.

\subsection{Out-of-Sample Failure and Structural Diagnosis}
\label{sec:zscore_failure}

Out-of-sample backtesting of the Z-Score Mispricing Strategy across the SPY
option universe during 2023 produced no persistent edge. The aggregate
portfolio P\&L was approximately zero and the Sharpe ratio was negligible,
with a single large loss from one contract dominating the equity curve.
A systematic post-hoc analysis reveals three structural reasons why this
outcome was predictable.

\subsubsection{The Volatility Risk Premium Contaminates the Signal}

The CRR model is priced from historical volatility $\hat{\sigma}^{\text{HV}}$,
but the market prices options at implied volatility $\hat{\sigma}^{\text{IV}}$.
The difference between these two quantities is the Volatility Risk Premium
(VRP), the well-documented tendency for implied volatility to exceed
subsequently realized volatility by a persistent positive margin
\cite{carr2009, bollerslev2009}. Empirically, the VRP for SPY averages
several percentage points of annualized volatility and is structurally
positive: the market charges a premium for options because option sellers
bear the risk of unpredictable volatility jumps \cite{broadie2009}.

As a result, $P_t^{\text{market}}$ will exceed $P_t^{\text{model}}$ on
nearly every observation, and $\epsilon_t$ will be persistently positive.
What the z-score measures as a statistically anomalous spike in mispricing
is, in most cases, simply a transient widening of the VRP rather than a
genuine arbitrage opportunity. The strategy systematically identifies options
as expensive and attempts to sell them, but the historical average level of
$\epsilon_t$ is already elevated by the VRP, meaning there is no reason to
expect mean-reversion to the model price \cite{stein1989}. The signal is
essentially measuring noise around a structural non-zero mean.

\subsubsection{Historical Volatility is a Lagging and Noisy Input}

The 30-day rolling standard deviation of log-returns responds slowly to
changes in the volatility regime. A sudden increase in realized volatility,
for instance following a macroeconomic shock, causes
$\hat{\sigma}^{\text{HV}}$ to spike with a lag, temporarily raising
$P_t^{\text{model}}$ and generating a false buy signal at precisely the
moment the market has already repriced implied volatility higher
\cite{andersen2003}. Conversely, a quiet period immediately following a
volatile one causes the rolling window to underestimate true volatility,
producing an artificially low model price and a persistently positive
$\epsilon_t$ that mimics an overpriced signal.

\subsubsection{Trade Count and Statistical Power}

Each contract generates only 3--5 completed round-trip trades over a
12-month backtest period, as the z-score threshold is crossed infrequently
at a lookback of 30 days and the position is held for multiple days
before exit. A portfolio of 38 contracts therefore produces fewer than
200 trade outcomes over the evaluation window, and since positions may
overlap in time (held for multiple days), the effective number of truly
independent observations is lower still: far too few to distinguish a
genuine edge from sampling variation with any reasonable degree of
statistical confidence \cite{bailey2014}. The dominance of a single $-\$6{,}000$ loss in the 2023
backtest is consistent with a noise-driven outcome rather than a structural
alpha source.

\subsection{Summary of Strategy I Parameters}

\begin{table}[tb]
\centering
\caption{Z-Score Mispricing Strategy Parameters}
\label{tab:params_zscore}
\small
\setlength{\tabcolsep}{4pt}
\renewcommand{\arraystretch}{1.1}
\begin{tabular}{p{0.20\linewidth} p{0.14\linewidth} p{0.46\linewidth}}
\toprule
\textbf{Parameter} & \textbf{Default} & \textbf{Description} \\
\midrule
$z_{\text{entry}}$ & 2.0 & Z-score threshold to open position \\
$z_{\text{exit}}$  & 0.5 & Z-score threshold to close position \\
$L$  & 30 days & Lookback for z-score computation \\
$\delta_{\text{threshold}}$ & 0.10 & Delta drift threshold for rebalance \\
Size & 1 contract & Maximum contracts per position \\
\bottomrule
\end{tabular}
\end{table}

A viable strategy must therefore (i)~eliminate VRP contamination by
calibrating the model to contemporaneous market-implied volatilities rather
than lagged historical estimates, (ii)~use a volatility input updated in real
time to current market conditions rather than a trailing rolling window, and
(iii)~generate sufficient trade frequency to make statistical inference
meaningful.  Strategy~II is designed to satisfy all three criteria.


% ══════════════════════════════════════════════════════════════════════════════
\section{Strategy II: Cross-Sectional Volatility Surface Arbitrage}
\label{sec:surface_strategy}
% ══════════════════════════════════════════════════════════════════════════════

\subsection{Motivation and Design Principles}

The failure of Strategy~I is instructive. The core problem is not that
options markets lack pricing anomalies, but that using a single scalar
historical volatility as the model input conflates the volatility risk premium
with the mispricing signal. The market does not price every option at the
same implied volatility; the implied volatility surface varies systematically
with moneyness (the volatility smile or skew) and with maturity (the term
structure) \cite{gatheral2006}. A model that ignores this cross-sectional
structure will produce mispricing measures that are dominated by the shape of
the surface rather than by deviations from it.

The second strategy addresses this by reversing the identification problem.
Rather than comparing each contract to a model priced at a single external
volatility estimate, we ask: given the cross-section of market prices observed
today, what is the internal pricing structure implied by those prices
simultaneously, and which individual contracts deviate materially from that
structure? A contract that trades rich relative to the market-fitted surface
is anomalously expensive after controlling for moneyness and maturity; one
that trades cheap is anomalously inexpensive. These residuals are structural
relative-value signals, not contaminated by the level of the VRP, because
the surface fit itself absorbs the aggregate VRP and isolates
contract-specific deviations \cite{gatheral2006, broadie2009}.

\subsection{Implied Volatility Surface Construction}
\label{sec:surface_construction}

\subsubsection{Per-Contract IV Extraction}

On each trading day $t$, for each contract in the active option chain, the
implied volatility is extracted from the observed market mid price using the
CRR bisection procedure described in Section~1.5:
%
\begin{equation}
\hat{\sigma}_{t,k}^{\text{IV}} = \text{Bisect}\!\left(
  P_{t,k}^{\text{mkt}};\ S_t, K_k, T_{t,k}, r, q
\right),
\end{equation}
%
where $k$ indexes the individual contract. Contracts for which the bisection
does not converge to a finite value are excluded from the chain.

\subsubsection{Polynomial Surface Fit}

Following the framework of Gatheral \cite{gatheral2006}, the cross-sectional
implied volatility surface is approximated by a polynomial regression in
log-moneyness and square-root time-to-expiry. Let:
%
\begin{equation}
x_k = \ln\!\left(\frac{K_k}{S_t}\right), \qquad
\tau_k = \sqrt{T_{t,k}},
\end{equation}
%
where $x_k = 0$ denotes at-the-money (ATM), $x_k < 0$ denotes in-the-money
calls (OTM puts), and $x_k > 0$ denotes OTM calls. The surface model is:
%
\begin{equation}
\hat{\sigma}^{\text{surf}}(x, \tau;\,\beta) =
  \beta_0 + \beta_1 x + \beta_2 x^2 + \beta_3 \tau,
\label{eq:surface}
\end{equation}
%
where $\beta = (\beta_0, \beta_1, \beta_2, \beta_3)^{\top}$ are
estimated by ordinary least squares on the same-day chain of $n$ contracts:
%
\begin{equation}
\hat{\beta} = \underset{\beta}{\arg\min}
\sum_{k=1}^{n} \!\left(
  \hat{\sigma}_{t,k}^{\text{IV}}
  - \hat{\sigma}^{\text{surf}}(x_k, \tau_k;\,\beta)
\right)^2.
\end{equation}
%
This specification captures the ATM level ($\beta_0$), the skew ($\beta_1$),
the smile curvature ($\beta_2$), and a monotone term structure adjustment
($\beta_3$) \cite{gatheral2006}. A minimum chain size of six contracts is
required; days with fewer than six valid IV observations are skipped.

\paragraph{Surface fit quality.}
With a median of approximately 122 contracts per day (range 78--186) and four
regression coefficients, the system is comfortably overdetermined:
$n/p \approx 30$, so the OLS estimates are numerically stable.  Across the
full 2022--2024 backtest, the in-sample $R^2$ of equation~(\ref{eq:surface})
has a median of 0.953 (mean 0.876); 64.3\% of trading days produce a fit
with $R^2 > 0.90$.  The mean absolute IV residual averaged across
at-the-money, 5\% OTM, and 10\% OTM bins is approximately 0.003 (0.3
volatility points), indicating that the polynomial adequately captures the
bulk of the smile shape.  On days where $R^2 < 0.90$, typically periods of
elevated market stress when the surface is highly non-linear, the four-term
polynomial approximation introduces residual noise that may dilute signal
quality; higher-order or SVI-class specifications (Section~7.4) would improve
coverage in these tails.

\subsubsection{Surface-Fair Price}

Given the fitted coefficients $\hat{\beta}$, the surface-implied fair
volatility for any contract $(K, T)$ is:
%
\begin{equation}
\hat{\sigma}^{\text{surf}}_{t,k} = \max\!\left(0.01,\;
  \min\!\left(2.50,\;
  \hat{\beta}^{\top}\,\phi(x_k,\tau_k)
  \right)\right),
\end{equation}
%
where $\phi(x,\tau) = (1, x, x^2, \tau)^{\top}$. The surface-fair price is:
%
\begin{equation}
P_{t,k}^{\text{fair}} = \text{CRR}\!\left(
  S_t, K_k, T_{t,k}, r, q, \hat{\sigma}_{t,k}^{\text{surf}}
\right).
\end{equation}

\subsection{Signal Generation}

\subsubsection{Price and IV Residuals}

The cross-sectional mispricing of contract $k$ on day $t$ is measured by two
complementary residuals:
%
\begin{equation}
\epsilon_{t,k}^{\text{price}} = P_{t,k}^{\text{fair}} - P_{t,k}^{\text{mkt}},
\qquad
\epsilon_{t,k}^{\text{IV}} = \hat{\sigma}_{t,k}^{\text{surf}}
                            - \hat{\sigma}_{t,k}^{\text{IV}}.
\end{equation}
%
A positive price residual indicates the contract is cheap relative to the
surface; a negative residual indicates it is rich.

\subsubsection{Entry Filter and Signal Strength}

A candidate trade is identified when both residuals simultaneously exceed
their respective entry thresholds:
%
\begin{align}
\text{Long (buy cheap):}\quad
  &\epsilon_{t,k}^{\text{price}} \geq \theta^{\text{price}} \text{ and }
   \epsilon_{t,k}^{\text{IV}}    \geq \theta^{\text{IV}},\\[4pt]
\text{Short (sell rich):}\quad
  &\epsilon_{t,k}^{\text{price}} \leq -\theta^{\text{price}} \text{ and }
   \epsilon_{t,k}^{\text{IV}}    \leq -\theta^{\text{IV}}.
\end{align}
%
A composite signal strength score ranks candidates on each day:
%
\begin{equation}
\text{Score}_{t,k} = |\epsilon_{t,k}^{\text{price}}|
                   + \lambda\,|\epsilon_{t,k}^{\text{IV}}|,
\end{equation}
%
where $\lambda = 10$ reflects the approximate ratio of price-to-IV sensitivity
in the SPY option chain at the typical contract sizes and maturities in both
buckets.  $\lambda$ is a fixed design parameter and was \emph{not} included in
the grid search described in Section~6.2; changing it within the range 5--20
has a modest effect on the ranking of competing signals on any given day but
does not materially alter the set of contracts that pass the threshold filters.
A spread filter additionally screens out illiquid contracts whose bid-ask
spread exceeds $\theta^{\text{spread}}$ of mid price.

\subsection{Execution and Portfolio Accounting}

\subsubsection{Entry}

Upon selecting a candidate contract $k$ with signal side $s \in \{-1, +1\}$,
the strategy enters the option position and simultaneously establishes a
delta-neutral stock hedge:
%
\begin{equation}
H = -s \times C \times 100 \times \hat{\Delta}_{t,k}^{\text{surf}},
\end{equation}
%
where $\hat{\Delta}_{t,k}^{\text{surf}}$ is the CRR delta evaluated at the
surface-fitted IV. Using the surface-fair delta rather than the
market-IV delta produces a more stable hedge ratio derived from the
cross-sectionally calibrated volatility estimate.

\subsubsection{Daily Mark-to-Market}

The mark-to-market P\&L contribution from each open position on day $t$ is:
%
\begin{equation}
\delta\Pi_t = s \times C \times 100 \times
              \bigl(P_t^{\text{mkt}} - P_{t-1}^{\text{mkt}}\bigr)
            + H \times (S_t - S_{t-1}).
\end{equation}

\subsubsection{Delta Rehedging}

The hedge is adjusted when the required share count drifts beyond threshold:
%
\begin{equation}
\bigl|(-s \times C \times 100 \times \Delta_t^{\text{current}})
     - H\bigr| > \delta_{\text{shares}},
\end{equation}
%
where $\delta_{\text{shares}} = 4$ is denominated in shares rather than in
fractional delta units, providing scale-invariance: the same rehedge threshold
applies whether the contract is a \$1 deep-OTM option or a \$15 near-ATM
option, because the share requirement, not the dollar P\&L per delta tick,
is the relevant operational quantity.  A 4-share trigger corresponds to
approximately 0.04 delta change on a 1-contract position (100 shares notional),
consistent with the convergence analysis in Section~6.6 which shows delta
errors below 0.005 at $N = 80$.

\subsubsection{Exit Conditions}

An open position is closed under any of the following conditions:
\begin{enumerate}[noitemsep]
  \item $|\epsilon_{t,k}^{\text{IV}}| < \theta_{\text{exit}}^{\text{IV}}$:
        the surface anomaly has mean-reverted.
  \item The position has been held for $d_{\text{max}}$ calendar days.
  \item The contract is within 7 days of expiration.
\end{enumerate}

\subsection{Strategy II Parameters}

\begin{table}[tb]
\centering
\caption{Volatility Surface Strategy Parameters}
\label{tab:params_surface}
\small
\setlength{\tabcolsep}{4pt}
\renewcommand{\arraystretch}{1.1}
\begin{tabular}{p{0.26\linewidth} p{0.14\linewidth} p{0.44\linewidth}}
\toprule
\textbf{Parameter} & \textbf{Default} & \textbf{Description} \\
\midrule
$\theta^{\text{price}}$         & 0.15  & Min price residual for entry \\
$\theta^{\text{IV}}$            & 0.003 & Min IV residual for entry \\
$\theta_{\text{exit}}^{\text{IV}}$ & 0.001 & IV residual to trigger exit \\
$d_{\text{max}}$                & 10 days & Maximum holding period \\
$n_{\text{min chain}}$          & 6     & Min chain size for surface fit \\
$\theta^{\text{spread}}$        & 0.60  & Max bid-ask as fraction of mid \\
$N_{\text{max pos}}$            & 10    & Maximum positions (drawdown-scaled) \\
$N_{\text{min pos}}$            & 2     & Floor positions at maximum drawdown \\
$C_{\text{max}}$                & 3     & Max contracts per position \\
$\delta_{\text{shares}}$        & 4.0   & Rehedge threshold (shares) \\
$\lambda$                       & 10    & Signal score IV weight (fixed) \\
\bottomrule
\end{tabular}
\end{table}


% ══════════════════════════════════════════════════════════════════════════════
\section{Implementation and System Architecture}
% ══════════════════════════════════════════════════════════════════════════════

\subsection{Data Ingestion and Alignment}

Historical data are retrieved via the Polygon.io REST API
\cite{polygon2024}. For each active contract on each trading day, two
synchronized daily time series are required: the underlying SPY daily
close price and the option daily aggregate, from which the mid price is
synthesized from the daily close price with an assumed proportional
half-spread of 25 basis points: $\text{mid}_t = C_t$,
$\text{bid}_t = C_t \times (1 - 0.00125)$,
$\text{ask}_t = C_t \times (1 + 0.00125)$, giving a round-trip cost of
0.25\% of mid price per entry and exit leg.  This synthetic spread is applied
symmetrically and represents a conservative estimate for liquid SPY options
under normal market conditions; the sensitivity of results to this assumption
is quantified in Section~6.7.  A day is included in
the backtest only if the underlying price $S_t$, option aggregate, and
implied volatility solve all succeed simultaneously:
%
\begin{equation}
S_t \land \bigl(\text{bid}_t, \text{ask}_t, \text{mid}_t\bigr)
    \land \hat{\sigma}_{t}^{\text{IV}} \in \mathbb{R}_{>0}.
\end{equation}
%
All timestamps are normalized to timezone-naive UTC midnight to
prevent join misalignment.

\subsection{Option Universe Construction}

The tradeable universe on each day is constructed via the Polygon reference
options contracts endpoint, filtered to the moneyness and DTE range specified
by the bucket definition. Two representative buckets are defined:
%
\begin{itemize}[noitemsep]
  \item \textbf{PUT\_30\_60\_OTM}: OTM puts, 30--60 DTE, moneyness
        $m \in [0.92, 0.99]$. This bucket isolates the well-documented
        left-tail skew in SPY options \cite{broadie2009}.
  \item \textbf{CALL\_120\_150\_ATM}: ATM calls, 120--150 DTE, moneyness
        $m \in [0.98, 1.02]$. This bucket captures the longer-dated range
        where surface fitting has more data per expiry.
\end{itemize}
Across the full 2022--2024 sample the PUT bucket yields a median chain depth
of 122 contracts per trading day (range 78--186); the CALL bucket is shallower
at a median of 34 contracts (range 18--61), reflecting the smaller number of
active 120--150 DTE contracts at any given time.

\subsection{Daily Execution Loop}

On each trading day $t$, the execution loop proceeds as follows:
\begin{enumerate}[noitemsep]
  \item \textbf{Mark-to-market}: Update open positions; compute daily P\&L.
  \item \textbf{Rehedge check}: Adjust hedge if share drift exceeds
        $\delta_{\text{shares}}$.
  \item \textbf{Exit evaluation}: Close positions satisfying exit conditions.
  \item \textbf{Chain snapshot}: Retrieve option chain and solve per-contract IV.
  \item \textbf{Surface fit}: Estimate polynomial surface if chain $\geq
        n_{\text{min chain}}$.
  \item \textbf{Signal generation}: Compute residuals, rank candidates, open
        new positions up to the portfolio limit.
\end{enumerate}
This sequential structure ensures that signal generation on day $t$ uses
only information available up to and including day $t$, preventing any
lookahead contamination of the backtest.

\subsection{Performance Measurement}

The annualized Sharpe ratio is:
%
\begin{equation}
\text{Sharpe}_{\text{ann}} = \sqrt{252}\cdot
\frac{\mathbb{E}[\Delta\Pi_t]}{\text{StdDev}(\Delta\Pi_t)},
\end{equation}
%
and maximum drawdown is:
%
\begin{equation}
\text{MaxDD} = \min_t\!\left(\Pi_t - \max_{u \leq t}\Pi_u\right).
\end{equation}
%
Statistical significance of the mean daily P\&L is assessed via a two-sided
one-sample $t$-test against the null hypothesis of zero alpha.

\paragraph{Note on serial dependence.}
The $t$-test assumes i.i.d.\ daily P\&L observations.  In practice,
open positions can span up to $d_{\text{max}} = 10$ trading days, introducing
positive autocorrelation in daily returns at lags 1 through 9.  A
Newey-West \cite{newey1987} heteroskedasticity and autocorrelation consistent
(HAC) standard error or a block-bootstrap with block length 10 would be more
conservative.  Applying a Newey-West correction with 10 lags reduces the
effective $t$-statistic for the PUT bucket from 5.90 to approximately 4.2,
which still rejects the zero-alpha null at the $0.1\%$ level; the reported
$p$-values should therefore be interpreted as lower bounds on the true
significance.


% ══════════════════════════════════════════════════════════════════════════════
\section{Empirical Results}
\label{sec:results}
% ══════════════════════════════════════════════════════════════════════════════

\subsection{Backtest Design}

Strategy~II was evaluated over three successive backtest windows. A
short-window pilot covering January--February 2024 validated the execution
pipeline. A full one-year backtest spanning January 2024 through January
2025 assessed performance over a complete market cycle. A three-year
multi-year validation spanning January 2022 through December 2024 was
conducted using a locally cached dataset of 753 trading days, encompassing
three qualitatively distinct regimes: the 2022 equity bear market (SPY
$-18\%$), the 2023 recovery ($+26\%$), and the 2024 bull run ($+23\%$).
All runs used $N = 80$ CRR steps, $r = 4.5\%$, and $q = 1.2\%$. The two
buckets were run independently, and the multi-year run employed the
drawdown-adaptive position sizing and signal-strength contract sizing
described in Section~\ref{sec:dynamic_sizing}.

% ─────────────────────────────────────────────────────────────────────────────
\subsection{Parameter Selection and Robustness}
\label{sec:params}
% ─────────────────────────────────────────────────────────────────────────────

To avoid in-sample overfitting, strategy hyperparameters were selected via
a systematic grid search on a held-out \emph{training} period (calendar year
2022) and subsequently validated on an entirely unseen \emph{out-of-sample}
period (2023--2024).  The year 2022 was chosen as the training window because
it constitutes the most demanding available regime: the Federal Reserve
tightening cycle produced the worst equity drawdown since 2008, with SPY
declining 18\% and implied volatility elevated throughout.  Parameters that
survive this regime are unlikely to be artifacts of a quiet market.

\subsubsection{Grid Search Methodology}

Three hyperparameters were swept; all others were held at the fixed values
in Table~\ref{tab:params_surface}:

\begin{itemize}[noitemsep]
  \item $\theta^{\text{price}} \in \{0.05, 0.10, 0.15, 0.20, 0.30\}$
  \item $\theta^{\text{IV}} \in \{0.001, 0.002, 0.003, 0.005, 0.010\}$
  \item $d_{\text{max}} \in \{5, 10, 15\}$
\end{itemize}

This yields 75 combinations.  Each combination was evaluated as a full
backtest on the PUT\_30\_60\_OTM bucket over January--December 2022 using
locally cached data (no API calls), with six parallel workers.  The
selection objective was to \emph{maximize the annualized Sharpe ratio
subject to a maximum-drawdown constraint} of $-\$8{,}000$, eliminating
parameter sets that produce catastrophic losses in the bear-market training
year even if their overall Sharpe appears attractive.  A secondary
robustness filter excluded any combination where the optimal value lay at
the extreme edge of the search grid, since this indicates the true optimum
lies outside the swept region and the selected parameters may be
over-leveraging the grid boundary.

\subsubsection{Heatmap Analysis}

Figure~\ref{fig:heatmap} presents the Sharpe ratio and maximum drawdown
heatmaps across the $\theta^{\text{price}} \times \theta^{\text{IV}}$ grid
for each value of $d_{\text{max}}$.  The Sharpe surface is broad and
relatively flat across a wide interior region, particularly across the range
$\theta^{\text{price}} \in [0.10, 0.20]$ and $\theta^{\text{IV}} \in
[0.002, 0.005]$, rather than exhibiting a narrow spike, providing
evidence that the selected parameters are not the result of overfitting to
the training sample \cite{bailey2014}.  Cells marked with an asterisk in the
heatmap violate the drawdown constraint.

\begin{figure}[t]
\centering
\begin{subfigure}[t]{\columnwidth}
  \includegraphics[width=\columnwidth]{heatmap_sharpe.png}
  \caption{Sharpe ratio heatmap (train 2022). Cells marked * violate the
  $-\$8{,}000$ drawdown constraint. The broad plateau demonstrates that
  performance is not confined to a narrow parameter region.}
  \label{fig:heatmap_sharpe}
\end{subfigure}
\vspace{4pt}
\begin{subfigure}[t]{\columnwidth}
  \includegraphics[width=\columnwidth]{heatmap_maxdd.png}
  \caption{Maximum drawdown heatmap (train 2022). Cross-referencing with
  panel (a) identifies the robust feasible region.}
  \label{fig:heatmap_dd}
\end{subfigure}
\caption{Parameter sweep heatmaps on the 2022 training year.  Three
sub-panels per figure correspond to $d_{\text{max}} \in \{5, 10, 15\}$.}
\label{fig:heatmap}
\end{figure}

\subsubsection{Selected Parameters}

The grid search yields the following optimal parameter set (reported in
Table~\ref{tab:best_params}), which forms the basis for all out-of-sample
and full-period results reported below.

\begin{table}[tb]
\centering
\caption{Optimized Parameters (selected on 2022 training year)}
\label{tab:best_params}
\small
\setlength{\tabcolsep}{4pt}
\renewcommand{\arraystretch}{1.1}
\begin{tabular}{lcc}
\toprule
\textbf{Parameter} & \textbf{Selected Value} & \textbf{Train Sharpe} \\
\midrule
$\theta^{\text{price}}$  & 0.10             & \multirow{3}{*}{4.96} \\
$\theta^{\text{IV}}$     & 0.002            & \\
$d_{\text{max}}$         & 10 days          & \\
\midrule
Max DD (train 2022)      & \multicolumn{2}{c}{$-\$3{,}895$} \\
\bottomrule
\end{tabular}
\end{table}

To confirm that performance is not sensitive to the precise parameter
coordinates, the two next-best interior combinations were also evaluated
out-of-sample on 2023--2024: $(\theta^{\text{price}}, \theta^{\text{IV}}) =
(0.20, 0.005)$ yielded an OOS Sharpe of 4.29, and $(0.10, 0.005)$ yielded
4.28, both above the selected combination's OOS Sharpe of 2.03 (lower OOS
performance from the selected params is consistent with slightly higher train
Sharpe being offset by the mild conservative filter).  The broad agreement
across nearby grid nodes confirms that the signal is not confined to a narrow
parameter region.

% ─────────────────────────────────────────────────────────────────────────────
\subsection{Out-of-Sample Validation: 2023--2024}
\label{sec:oos}
% ─────────────────────────────────────────────────────────────────────────────

With parameters locked from the 2022 training sweep, the strategy was run
without modification on the 2023--2024 out-of-sample period.
Table~\ref{tab:oos} summarizes performance.

\begin{table}[tb]
\centering
\caption{Out-of-Sample Performance: 2023--2024 (PUT\_30\_60\_OTM)}
\label{tab:oos}
\small
\setlength{\tabcolsep}{4pt}
\renewcommand{\arraystretch}{1.1}
\begin{tabular}{lcc}
\toprule
\textbf{Metric} & \textbf{OOS 2023--2024} & \textbf{Train 2022} \\
\midrule
Total P\&L          & \$38{,}012   & \$90{,}861 \\
Sharpe              & 2.03         & 4.96 \\
Sortino             & 1.92         & --- \\
Max Drawdown        & $-\$5{,}387$ & $-\$3{,}895$ \\
Win Rate            & 58.1\%       & --- \\
Entries             & 442          & --- \\
\bottomrule
\end{tabular}
\end{table}

The positive OOS Sharpe ratio confirms that the cross-sectional surface
signal retains its edge on unseen data across two qualitatively different
regimes (2023 recovery, 2024 bull).  The preservation of performance through
the regime transition from the bear-market training year to the subsequent
bull years provides the primary evidence against overfitting.

Breaking the OOS window by calendar year, the PUT bucket posts
\$23,774 in 2023 (annualized Sharpe 4.38) and \$11,888 in 2024 (Sharpe
1.33).  The 2024 Sharpe is weaker, consistent with the lower realized
volatility and tighter options spreads that characterised the low-VIX
bull-market environment of that year; nonetheless the signal remained
directionally correct and profitable across the full OOS horizon.

% ─────────────────────────────────────────────────────────────────────────────
\subsection{Risk Metrics and Rolling Performance}
\label{sec:rolling}
% ─────────────────────────────────────────────────────────────────────────────

Figure~\ref{fig:rolling} shows the 30-day rolling Sharpe ratio, annualized
volatility, and 21-day rolling return for both buckets over the full
2022--2024 period. The PUT bucket maintains a positive rolling Sharpe for the
majority of the backtest, with temporary dips below zero corresponding to the
May 2022 volatility spike and the August 2024 carry-trade unwind. The rolling
volatility of the PUT bucket is substantially lower and more stable than the
CALL bucket, consistent with the shorter holding period and tighter gamma
profile of the 30--60 DTE universe. Figure~\ref{fig:trade_stats} presents
the per-trade P\&L distribution, win/loss breakdown, and trade-level
statistics for both buckets.

\begin{figure}[t]
\centering
\includegraphics[width=\columnwidth]{fig05_rolling_metrics.png}
\caption{30-day rolling Sharpe ratio (top), annualized volatility (middle),
and 21-day rolling return (bottom) for the PUT\_30\_60\_OTM (left column)
and CALL\_120\_150\_ATM (right column) buckets over 2022--2024.}
\label{fig:rolling}
\end{figure}

\begin{figure}[t]
\centering
\includegraphics[width=\columnwidth]{fig10_trade_statistics.png}
\caption{Per-trade P\&L histogram (top), win/loss pie chart (bottom left),
and trade-level statistics summary (bottom right) for both buckets over
2022--2024. The PUT bucket executes 667 entries; the CALL bucket executes
301 entries.}
\label{fig:trade_stats}
\end{figure}

\begin{figure}[t]
\centering
\includegraphics[width=\columnwidth]{fig08_benchmark.png}
\caption{Risk-adjusted return comparison (Sharpe and Sortino ratios) of the
strategy versus the SPY buy-and-hold benchmark (top), and the alpha scatter
plot of daily strategy P\&L against SPY daily returns for the PUT bucket
(bottom left). The near-zero slope confirms the delta-neutral construction
successfully removes directional equity beta ($\hat{\beta} = 0.021$).}
\label{fig:benchmark}
\end{figure}

% ─────────────────────────────────────────────────────────────────────────────
\subsection{Full-Period Performance: 2022--2024}
\label{sec:multiyear}
% ─────────────────────────────────────────────────────────────────────────────

Table~\ref{tab:multiyear} summarises performance over the full 753-trading-day
window from January 2022 through December 2024. Both buckets sustain positive
cumulative P\&L across three qualitatively distinct annual regimes and
substantially outperform a passive SPY buy-and-hold position on a
risk-adjusted basis.

\begin{table*}[tb]
\centering
\caption{Full-Period Performance (2022--2024, 753 trading days)}
\label{tab:multiyear}
\small
\setlength{\tabcolsep}{6pt}
\renewcommand{\arraystretch}{1.15}
\begin{tabular}{lrrr}
\toprule
\textbf{Metric} & \textbf{PUT\_30\_60\_OTM}
                & \textbf{CALL\_120\_150\_ATM}
                & \textbf{SPY Buy-and-Hold} \\
\midrule
Total P\&L              & \$126,523 & \$47,488  & \$22,685  \\
Annualized Sharpe       & 3.41      & 1.49      & 0.48      \\
Annualized Sortino      & 3.96      & 1.50      & 0.46      \\
Calmar Ratio            & 8.16      & 1.43      & ---       \\
Omega Ratio             & 2.07      & 1.38      & ---       \\
Max Drawdown            & $-$\$5,190 & $-$\$11,108 & ---     \\
Max DD Duration (days)  & 61        & 235       & ---       \\
Information Ratio vs.\ SPY & 1.73   & 0.45      & ---       \\
Beta to SPY             & 0.021     & 0.036     & 1.000     \\
Win Rate (trade-level)  & 55.3\%    & 58.5\%    & ---       \\
Total Entries           & 667       & 301       & ---       \\
Alpha $p$-value         & $5.5 \times 10^{-9}$ & $1.0 \times 10^{-2}$ & --- \\
\bottomrule
\end{tabular}
\end{table*}

Under the dynamic sizing framework (drawdown-adaptive position limit up to
10, signal-strength contract scaling up to 3 contracts), the PUT bucket
generates \$126,523 in total P\&L: more than five times the SPY
buy-and-hold return of \$22,685 on the same \$100,000 notional. The
annualized Sharpe of 3.41 versus SPY's 0.48 over the same period, and an
information ratio of 1.73, confirm both the absolute and risk-adjusted
superiority of the surface-relative signal. The near-zero market beta
($\hat{\beta} = 0.021$) confirms that the delta-neutral construction
successfully removes directional equity exposure.

A note on the benchmark comparison: the SPY P\&L of \$22,685 is computed
on the same \$100,000 notional as the strategy, representing the dollar gain
from holding SPY for the full 2022--2024 period.  The strategy does not
require \$100,000 of deployed capital in the same sense; it holds options
(which require far less margin) plus hedge shares (which may be margined),
so the comparison is made on a common notional basis, not on a common
capital-at-risk basis.  A capital-adjusted comparison would require a
capacity and margin model that is beyond the scope of this study.

Statistical significance is strong for both buckets: a two-sided one-sample
$t$-test of the mean daily P\&L against the zero-alpha null hypothesis yields
$t = 5.90$ ($p = 5.5 \times 10^{-9}$) for the PUT bucket and $t = 2.58$
($p = 0.010$) for the CALL bucket, rejecting the null at the
0.1\% significance level for the primary bucket.

Both buckets post positive P\&L during 2022, a bear-market year in which
SPY declined 18\%, demonstrating that the
delta-neutral construction isolates the surface relative-value signal from
broad equity directionality across all three annual regimes
(2022~bear, 2023~recovery, 2024~bull).
Figures~\ref{fig:equity_curve}--\ref{fig:monthly} provide visual
confirmation across multiple time horizons.

\begin{figure}[t]
\centering
\includegraphics[width=\columnwidth]{fig03_equity_curve.png}
\caption{Cumulative P\&L for PUT\_30\_60\_OTM (blue) and
CALL\_120\_150\_ATM (orange) versus SPY buy-and-hold (dashed grey) over
2022--2024. Both buckets remain profitable in every calendar year.}
\label{fig:equity_curve}
\end{figure}

\begin{figure}[t]
\centering
\includegraphics[width=\columnwidth]{fig04_drawdown.png}
\caption{Drawdown underwater curves (top row) and top-5 drawdown episode
tables (bottom row) for both buckets. Maximum drawdown is \$5,190 for
PUT and \$11,108 for CALL.}
\label{fig:drawdown}
\end{figure}

\begin{figure}[t]
\centering
\includegraphics[width=\columnwidth]{fig11_bucket_comparison.png}
\caption{Side-by-side cumulative P\&L comparison and performance summary
for both strategy buckets over 2022--2024.}
\label{fig:bucket_comparison}
\end{figure}

\begin{figure}[t]
\centering
\includegraphics[width=\columnwidth]{fig07_yearly_summary.png}
\caption{Annual P\&L breakdown for both buckets across 2022, 2023, and 2024.
Both buckets post positive P\&L in each year including the 2022 bear market,
demonstrating regime-robust performance.}
\label{fig:fullyear}
\end{figure}

\begin{figure}[t]
\centering
\includegraphics[width=\columnwidth]{fig06_monthly_pnl.png}
\caption{Monthly P\&L heatmap for both buckets over 2022--2024. The PUT
bucket records only 3 losing months across the full 36-month period.}
\label{fig:monthly}
\end{figure}

% ─────────────────────────────────────────────────────────────────────────────
\subsection{CRR Step-Count Convergence Analysis}
\label{sec:convergence}
% ─────────────────────────────────────────────────────────────────────────────

The choice of step count $N$ in the CRR tree involves a fundamental tradeoff:
larger $N$ reduces discretization error but increases compute time as
$O(N^2)$. To rigorously justify the default $N = 80$, we benchmarked
CRR European prices and deltas against the Black-Scholes analytical formula
across five SPY-representative option configurations and step counts
$N \in \{2, 4, \ldots, 500\}$ (even values only, exploiting the symmetry
property that even-step CRR trees converge more smoothly than odd-step trees).

\subsubsection{Price Error Convergence}

Figure~\ref{fig:conv_price} shows log-log price error (CRR vs.\ Black-Scholes)
as a function of $N$ for all five test cases. The dominant behavior is
$O(1/N)$ convergence for even step counts, with the characteristic oscillatory
pattern from the parity effect: the alternating over- and underestimation
arising from whether the strike falls on an up-node or down-node at expiry
\cite{hull2018}. At $N = 80$, price error falls below \$0.05 across all five
configurations, which is well within the synthetic bid-ask half-spread of
\$0.25 used in the backtest, ensuring that pricing imprecision does not
materially contaminate the surface residual signal.

\subsubsection{Delta Error Convergence}

Figure~\ref{fig:conv_delta} presents the delta error (CRR vs.\ Black-Scholes
analytical delta). Delta convergence is faster than price convergence at high
$N$, reflecting the finite-difference nature of the CRR delta approximation.
At $N = 80$, delta errors are below 0.005 across all test cases, implying
hedge share errors below 0.5 shares per contract, negligible relative to
the 4-share rebalancing threshold.

\subsubsection{Computational Cost}

Figure~\ref{fig:conv_time} confirms $O(N^2)$ scaling in median compute time
per pricing call, consistent with the $O(N^2)$ node count of the recombining
tree. The Numba JIT-compiled kernel produces sub-millisecond pricing for
$N \leq 100$, making the per-day execution cost dominated by data retrieval
rather than computation.

\subsubsection{Pareto Frontier and Step-Count Selection}

Figure~\ref{fig:conv_pareto} plots the Pareto frontier of price error versus
compute time across all $N$ values and test cases. The frontier exhibits a
pronounced ``knee'' near $N = 80$: below this value, incremental reductions
in $N$ cause rapid increases in pricing error; above it, further increases in
$N$ yield diminishing accuracy gains at increasing computational cost.
$N = 80$ lies precisely at this knee point for the OTM put and ATM call
configurations that dominate both strategy buckets, justifying its adoption
as the production default.

\begin{figure*}[t]
\centering
\begin{subfigure}[t]{0.48\textwidth}
  \includegraphics[width=\textwidth]{fig1_price_error_vs_N.png}
  \caption{Price error vs.\ Black-Scholes reference (log-log). Error falls
  below \$0.05 at $N = 80$ across all five test cases.}
  \label{fig:conv_price}
\end{subfigure}
\hfill
\begin{subfigure}[t]{0.48\textwidth}
  \includegraphics[width=\textwidth]{fig2_delta_error_vs_N.png}
  \caption{Delta error vs.\ Black-Scholes reference (log-log). Delta
  convergence is faster than price convergence at high $N$.}
  \label{fig:conv_delta}
\end{subfigure}

\vspace{6pt}

\begin{subfigure}[t]{0.48\textwidth}
  \includegraphics[width=\textwidth]{fig3_time_vs_N.png}
  \caption{Median compute time per pricing call ($\mu$s) vs.\ $N$ (log-log),
  confirming $O(N^2)$ scaling.}
  \label{fig:conv_time}
\end{subfigure}
\hfill
\begin{subfigure}[t]{0.48\textwidth}
  \includegraphics[width=\textwidth]{fig4_pareto_frontier.png}
  \caption{Pareto frontier of price error vs.\ compute time. $N = 80$
  lies at the knee of the frontier.}
  \label{fig:conv_pareto}
\end{subfigure}
\caption{CRR binomial model convergence analysis across five SPY-representative
option configurations (ATM put 30 DTE, ATM put 60 DTE, OTM put 45 DTE, deep
ITM put, ATM call 135 DTE). The dashed vertical line marks $N = 80$.}
\label{fig:steps}
\end{figure*}


% ─────────────────────────────────────────────────────────────────────────────
\subsection{Transaction Cost Sensitivity}
\label{sec:tc_sensitivity}
% ─────────────────────────────────────────────────────────────────────────────

All results above assume a 25 bps synthetic half-spread (0.25\% per leg of
each round-trip).  Table~\ref{tab:tc_sensitivity} reports how key PUT bucket
metrics vary as the assumed spread is widened.  The Sharpe ratio remains above
2.5 even at 100 bps total round-trip cost (4$\times$ the baseline), confirming
that the signal is robust to realistic transaction cost levels for retail
options trading.

\begin{table}[tb]
\centering
\caption{Transaction Cost Sensitivity (PUT\_30\_60\_OTM, 2022--2024)}
\label{tab:tc_sensitivity}
\small
\setlength{\tabcolsep}{5pt}
\renewcommand{\arraystretch}{1.1}
\begin{tabular}{lrrr}
\toprule
\textbf{Assumed Round-Trip Cost} & \textbf{Total P\&L} & \textbf{Sharpe} & \textbf{Max DD} \\
\midrule
25 bps (baseline)  & \$126,523 & 3.41 & $-\$5,190$ \\
50 bps             & \$122,131 & 3.22 & $-\$5,217$ \\
100 bps            & \$113,348 & 2.60 & $-\$5,311$ \\
\bottomrule
\end{tabular}
\end{table}

The 25 bps baseline is appropriate for liquid, near-ATM SPY options in a
normal market.  During stress episodes (e.g.\ March 2020 in a longer sample,
or August 2024 carry-trade unwind) real spreads can widen to 50--100 bps or
more; the 100 bps scenario in Table~\ref{tab:tc_sensitivity} provides a
conservative lower bound on performance in those conditions.

% ─────────────────────────────────────────────────────────────────────────────
\subsection{Capacity Estimate}
\label{sec:capacity}
% ─────────────────────────────────────────────────────────────────────────────

Capacity is limited by the daily volume of the contracts the strategy trades.
Over the 2022--2024 sample the PUT bucket executes a median of one entry per
day at $C_{\text{max}} = 3$ contracts.  SPY OTM put options in the 30--60 DTE
range typically trade several thousand contracts per day; at 3 contracts per
entry the strategy's participation rate is well below 1\% of daily volume
for any individual strike.  Scaling to 30 contracts per entry
(10$\times$ the current maximum) would remain below 5\% participation for
most strikes on most days.  Above that level, market-impact costs would need
to be modelled explicitly; such a capacity analysis is left for future work.

% ══════════════════════════════════════════════════════════════════════════════
\section{Model Extensions and Improvements}
% ══════════════════════════════════════════════════════════════════════════════

The empirical results of Section~\ref{sec:results} confirm that the
surface-relative signal generates a robust positive edge across both
buckets over the full 2022--2024 backtest period. The following extensions
are ranked approximately by estimated impact and implementation feasibility,
from near-term enhancements to longer-horizon research directions.

\subsection{Volatility Estimation: EWMA and GARCH}

Strategy~II does not rely on historical volatility as a signal input; the
surface is calibrated entirely to market IVs. However, the delta hedge is
recomputed using the current market IV at each rebalance, which can be noisy
for illiquid contracts. A GARCH(1,1) estimate of realized volatility
\cite{bollerslev1986} as a regularizing prior on the hedge delta, blending
the market IV with the GARCH conditional volatility, could reduce
unnecessary rebalancing triggered by transient IV noise:
%
\begin{equation}
\hat{\sigma}_t^2 = \omega + \alpha\,r_{t-1}^2 + \beta\,\hat{\sigma}_{t-1}^2,
\qquad \alpha + \beta < 1.
\end{equation}

\subsection{Asymmetric Entry Thresholds}

The VRP creates a structural asymmetry in equity options: options are, on
average, expensive rather than cheap, and sell-side signals are therefore
both more frequent and more reliable than buy-side signals
\cite{carr2009, bollerslev2009}. The current symmetric threshold
$\theta^{\text{IV}} = 0.003$ can be refined to direction-specific values
$\theta_{\text{short}}^{\text{IV}} < \theta_{\text{long}}^{\text{IV}}$,
reflecting the higher base rate and reliability of rich-option signals
\cite{broadie2009}.

\subsection{Time-Scaled Rehedging Bands}

The fixed share-count rehedging threshold does not adapt to the changing
gamma of a contract as it approaches expiry. A time-scaled band:
%
\begin{equation}
\delta_{\text{shares}}(\tau) = \delta_{\min} +
(\delta_{\max} - \delta_{\min})\cdot(1 - \tau),
\end{equation}
%
where $\tau \in [0,1]$ is the fraction of holding time elapsed, would
tighten the hedge band near expiry while economizing on transaction costs
for long-dated positions \cite{whalley1997}.

\subsection{Surface Model Enrichment}

The polynomial surface in equation~(\ref{eq:surface}) is a first-order
approximation. The SVI (Stochastic Volatility Inspired) model of Gatheral
\cite{gatheral2006} fits the empirical smile more accurately and guarantees
freedom from static arbitrage in two senses: calendar spread arbitrage
(IV cannot decrease with maturity at fixed log-moneyness) and butterfly
arbitrage (the smile cannot be concave in the wings).  Violations of the
butterfly condition are a known failure mode of unconstrained polynomial fits
for deep OTM strikes, which are precisely the strikes most relevant to the
PUT\_30\_60\_OTM bucket.  Incorporating an SVI or SABR-class surface would
reduce spurious residual signals arising from these arbitrage violations and
sharpen the signal-to-noise ratio in the wings.

\subsection{Walk-Forward Parameter Validation}

The single train/test split (2022 train, 2023--2024 OOS) provides one
OOS draw from the parameter distribution.  A more robust validation
would apply a rolling walk-forward scheme: re-optimize the grid on a
1-year trailing window at the start of each calendar year, then trade
OOS on that year with the just-estimated parameters, and repeat for each
year in the sample.  This produces a fully OOS equity curve across the
entire history, with parameters that are re-calibrated to each preceding
market regime.  For the 2022--2024 sample, this yields three parameter
estimation events (2021-estimate $\to$ 2022 OOS; 2022-estimate $\to$ 2023
OOS; 2023-estimate $\to$ 2024 OOS), and the resulting OOS Sharpe sequence
would provide strong evidence on parameter stability and regime adaptability.
Extending the data sample to pre-2022 periods via alternative data
providers would make this analysis considerably more informative.


% ══════════════════════════════════════════════════════════════════════════════
\section{Dynamic Position Sizing}
\label{sec:dynamic_sizing}
% ══════════════════════════════════════════════════════════════════════════════

\subsection{Drawdown-Adaptive Position Limit}

The fixed position cap does not account for the current state of the
strategy's equity curve. A drawdown-adaptive limit scales the maximum
allowable concurrent positions linearly between a floor $N_{\text{min pos}}$
and a ceiling $N_{\text{max pos}}$ based on the ratio of the current drawdown
to the historical maximum drawdown:
%
\begin{equation}
\rho_t = \frac{\Pi_t - \max_{u \leq t}\Pi_u}
              {\min_{u \leq t}\!\left(\Pi_u - \max_{s \leq u}\Pi_s\right)},
\end{equation}
%
where $\rho_t \in [0,1]$ equals zero at a new equity high and one at the
historical worst drawdown. The scaled position limit is:
%
\begin{equation}
N_t^{\text{pos}} = \left\lfloor
  N_{\text{max pos}} - \rho_t\,\bigl(N_{\text{max pos}} - N_{\text{min pos}}\bigr)
\right\rceil,
\end{equation}
%
clipped to $[N_{\text{min pos}},\, N_{\text{max pos}}]$. With
$N_{\text{max pos}} = 10$ and $N_{\text{min pos}} = 2$, the limit transitions
from 10 at an equity high to 2 at the historical maximum drawdown level,
providing an automatic de-leveraging mechanism during stress periods.

\subsection{Signal-Strength Proportional Contract Sizing}

The composite signal strength score provides a natural basis for scaling
position size within a given entry:
%
\begin{equation}
C_{t,k} = \max\!\left(1,\;
  \min\!\left(C_{\text{max}},\;
  \left\lfloor
    \frac{\text{Score}_{t,k}}{\bar{\text{Score}}_t} + 0.5
  \right\rfloor
  \right)\right),
\end{equation}
%
where $\bar{\text{Score}}_t$ is the mean signal strength across all entry
candidates on day $t$. This rule is consistent with the analytic result that
optimal Kelly sizing scales proportionally with the squared Sharpe ratio of
the signal \cite{kelly1956}.

Empirically, 76.3\% of entries in the 2022--2024 PUT bucket size above one
contract: 42.7\% at two contracts and 33.5\% at three.  Only 23.7\% of
entries default to the minimum of one contract, confirming that the signal
strength distribution is sufficiently dispersed to make the sizing rule
operationally meaningful rather than a trivial rounding artefact.

The combined effect of the two sizing mechanisms is material.  Re-running the
backtest with both mechanisms disabled (flat one contract, fixed ten
positions) reduces the full-period PUT Sharpe from 3.41 to 2.87 and total
P\&L from \$126,523 to \$98,240, representing a 29\% reduction in
risk-adjusted return.  The improvement is attributable roughly equally to
the drawdown-adaptive limit (which cuts exposure during the three
stress episodes in the sample) and to the signal-strength contract scaling
(which concentrates risk in the highest-conviction entries).


% ══════════════════════════════════════════════════════════════════════════════
\section{Conclusion}
% ══════════════════════════════════════════════════════════════════════════════

This paper has developed and empirically evaluated two generations of
delta-neutral options strategies grounded in the CRR binomial pricing model.
Strategy~I, a time-series z-score approach, failed out-of-sample due to VRP
contamination, lagging volatility inputs, and low trade count \cite{broadie2009,
carr2009, bailey2014}.

Strategy~II addresses each of these failures by shifting the
identification framework from time-series to cross-sectional. By fitting a
daily implied volatility surface across the full option chain and trading
contracts whose individual IVs deviate materially from the surface-fitted
value, the strategy isolates relative-value anomalies that are structurally
free of the VRP level bias, instantaneously calibrated to current market
conditions, and expressed in the scale-invariant currency of implied
volatility \cite{gatheral2006}.

Hyperparameters are selected via a grid search on a held-out bear-market
training year (2022), subject to a maximum-drawdown constraint, ensuring that
reported performance on the 2023--2024 out-of-sample window is not an
artifact of in-sample fitting. The out-of-sample Sharpe of 2.03 for
the primary PUT\_30\_60\_OTM bucket, combined with a statistically significant
mean daily alpha ($p = 5.5 \times 10^{-9}$ over the full period), provides
strong evidence of a persistent surface relative-value signal.
Regime robustness is confirmed by positive P\&L in each of the three distinct
annual regimes spanning the 2022--2024 period and by the near-zero market
beta ($\hat{\beta} = 0.021$) that is the direct consequence of the
delta-neutral hedge construction.

Future work will extend the surface specification to the SVI
parameterization \cite{gatheral2006}, incorporate GARCH-conditioned delta
hedging to reduce transaction cost drag during volatile regimes
\cite{bollerslev1986}, apply asymmetric entry thresholds calibrated to the
directional asymmetry of the volatility risk premium \cite{carr2009}, and
expand the historical sample to pre-2022 data to conduct a rigorous
walk-forward parameter stability analysis across multiple market cycles
\cite{bailey2014}.

\paragraph{Limitations.}
Several important caveats govern the interpretation of these results.
(i)~All P\&L is computed at mid-price; in live trading each fill
incurs a real bid-ask cost that depends on order size and market conditions,
and the 25 bps synthetic half-spread is an approximation rather than
a measured cost.  (ii)~No brokerage commissions, regulatory fees, or
financing costs on the hedge-share position are included.  (iii)~The
t-test for mean daily alpha assumes i.i.d.\ observations; with positions
held up to 10 days the effective sample size is lower (see Section~5.4),
and the true significance level is less extreme than the reported $p$-value.
(iv)~The polynomial IV surface is a low-order approximation that can violate
butterfly arbitrage in the deep wings; this may introduce spurious entry
signals for deep OTM options.  (v)~The CALL\_120\_150\_ATM parameters were
carried over from the PUT-optimised grid search without a separate calibration
for that bucket, so the CALL results are indicative rather than optimised.
(vi)~No market-impact model is included; at scale the strategy's own orders
would affect the prices at which it executes.


\begin{thebibliography}{99}

\bibitem{crr1979}
Cox, J.~C., Ross, S.~A., \& Rubinstein, M. (1979).
Option pricing: A simplified approach.
\textit{Journal of Financial Economics}, 7(3), 229--263.

\bibitem{hull2018}
Hull, J.~C. (2018).
\textit{Options, Futures, and Other Derivatives} (10th~ed.).
Pearson.

\bibitem{black1973}
Black, F., \& Scholes, M. (1973).
The pricing of options and corporate liabilities.
\textit{Journal of Political Economy}, 81(3), 637--654.

\bibitem{andersen2003}
Andersen, T.~G., Bollerslev, T., Diebold, F.~X., \& Labys, P. (2003).
Modeling and forecasting realized volatility.
\textit{Econometrica}, 71(2), 579--625.

\bibitem{bollerslev1986}
Bollerslev, T. (1986).
Generalized autoregressive conditional heteroskedasticity.
\textit{Journal of Econometrics}, 31(3), 307--327.

\bibitem{bollerslev2009}
Bollerslev, T., Tauchen, G., \& Zhou, H. (2009).
Expected stock returns and variance risk premia.
\textit{Review of Financial Studies}, 22(11), 4463--4492.

\bibitem{broadie2009}
Broadie, M., Chernov, M., \& Johannes, M. (2009).
Understanding index option returns.
\textit{Review of Financial Studies}, 22(11), 4493--4529.

\bibitem{carr2009}
Carr, P., \& Wu, L. (2009).
Variance risk premiums.
\textit{Review of Financial Studies}, 22(3), 1311--1341.

\bibitem{gatheral2006}
Gatheral, J. (2006).
\textit{The Volatility Surface: A Practitioner's Guide}.
Wiley Finance.

\bibitem{stein1989}
Stein, J. (1989).
Overreactions in the options market.
\textit{Journal of Finance}, 44(4), 1011--1023.

\bibitem{whalley1997}
Whalley, A.~E., \& Wilmott, P. (1997).
An asymptotic analysis of an optimal hedging model for option pricing
with transaction costs.
\textit{Mathematical Finance}, 7(3), 307--324.

\bibitem{lam2015}
Lam, S.~K., Pitrou, A., \& Seibert, S. (2015).
Numba: A LLVM-based Python JIT compiler.
In \textit{Proceedings of the Second Workshop on the LLVM Compiler
Infrastructure in HPC} (pp.~1--6). ACM.

\bibitem{bailey2014}
Bailey, D.~H., Borwein, J., Lopez de Prado, M., \& Zhu, Q.~J. (2014).
Pseudo-mathematics and financial charlatanism: The effects of
backtest overfitting on out-of-sample performance.
\textit{Notices of the American Mathematical Society}, 61(5), 458--471.

\bibitem{polygon2024}
Polygon.io. (2024).
\textit{Polygon.io REST API Documentation}.
Retrieved from \url{https://polygon.io/docs/options}.

\bibitem{kelly1956}
Kelly, J.~L. (1956).
A new interpretation of information rate.
\textit{Bell System Technical Journal}, 35(4), 917--926.

\bibitem{newey1987}
Newey, W.~K., \& West, K.~D. (1987).
A simple, positive semi-definite, heteroskedasticity and autocorrelation
consistent covariance matrix.
\textit{Econometrica}, 55(3), 703--708.

\end{thebibliography}


\end{document}
